\section{受迫与强导引场算例：需要逐篇保存的实际输入}
\subsection{Taylor问题的半隐式计算}
两篇Taylor研究均用$256\times128$网格、$\Delta t=10^{-3}$，$x$方向在电流层附近加密，空间与时间二阶，使用Harned--Miki\'c型半隐式处理。$x=\pm1$导体墙、$y$周期长度8，驱动上升时间1；常用超电阻$\nu=2.5\times10^{-9}$。任意导引场篇另取黏性和热扩散$\mu=\kappa=10^{-7}$。这些量按各文的系统尺度归一化，不能与F3D按$d_i$归一化的耗散系数直接比较。\citep{L057,L058}

两文给出了耗散与驱动扫描，但没有独立的网格、步长收敛系列。复算应先测X点电阻和超电阻的贡献，再在固定物理系数下细化；若同时降低超电阻，会把模型变化与离散误差混在一起。墙位移、墙速度与磁通边界必须彼此一致，具体关系已在第一章推导。

\subsection{F3D受迫算例：电子惯性不替代X点耗散账目}
2005年二维算例用$102.4d_i$方盒、$512\times1024$网格，$\Delta x=0.2d_i,\Delta y=0.1d_i$、质量比25；质量比100对照用$1024\times2048$。空间四阶差分，时间二阶trapezoidal leapfrog；密度和动量加四阶耗散$5.1\times10^{-5}$。2008年三维改为$25.6d_i$周期立方体、$128\times256\times128$，仍只有约两格跨法向$d_e$。\citep{L076,F717}

模型用常系数$\B^*=(1-(m_e/m_i)\nabla^2)\B$，并非一般变密度电子惯性。准稳态X点仍依赖数值扩散；所以“没有显式电阻”不能写成“X点非理想项全由惯性支撑”。三维没有给出充分分辨率扫描，第三维较粗也限制不稳定性研究。

复现前还需处理原文输入疑点。2008年所获版本的驱动时间函数印作$\tan$，与稳定渐开驱动的叙述及2005年的$\tanh$形式矛盾；$\tau_f$也有前后不同写法。这些应向输入文件追溯，不能用猜测的公式宣称逐图复现。式(17a)、(18a)的出流标度疑点则可由连续性与动量平衡独立发现。

\subsection{PIXIE2D与VPIC：局地对照和尺寸扫描用了不同参数}
\citet{L217}使用Newton--Krylov、物理多重网格预条件和向X点集中的网格；空间二阶中心差分，对流项用QUICK，时间可选Crank--Nicolson或BDF2。论文未逐算例列出格点数、最小网格和时间步，因此只能登记其方法与收敛陈述，不能补造完整输入。

\begin{center}\small
\begin{tabular}{lrrr}\toprule
VPIC比较 & 无快速色散波局地例 & 有快速色散波例 & 无快速色散波尺寸例\\\midrule
$m_i/m_e$ &15.2&380.3&15.2\\
$B_g/B_{\rm rec}$ &17.4&12&17.4\\
$v_{the}/c$ &0.0625&0.308&0.25\\
$\omega_{pe}/\Omega_{ce,x}$ &10&10&2.5\\
$T_i/T_e$ &0.25&0.1&0.25\\\bottomrule
\end{tabular}\end{center}
尺寸例的Alfvén归一化还考虑相对论修正。长文图14与15的流体黏性分别为$10^{-4}$与$10^{-5}$；有快速色散波的PIC例每格16000粒子。这些细节说明，不能把所有图拼成同一质量比、同一温度和同一耗散的单变量扫描。\citep{L216,L217}

\section{陀螺流体和非局地热流：比网格大小更容易漏掉的设置}
\subsection{四场双重加速的伪谱计算}
\citet{F672}把$0.1\operatorname{sech}^2x$平衡截成23个Fourier模，采用双周期边界、三阶Adams--Bashforth和伪谱求导，最高$4096\times512$，并在微尺度滤波。非线性系列$L_y=4\pi$、$k_y=0.5$、$\Delta'\simeq59.9$，$d_e=0.2$至0.025，同时保持$d_i=10d_e,\rho_s=2d_e,r_i=4d_e$。

所报告的一次$x$分辨率加倍，峰值率基本不变、峰值电子流改变不到6\%；它支持该例宏观率的稳健性，不等于所有细丝、所有参数均收敛。微层继续变细到不能解析时应停止判读。原文线性图2在正文写$\rho_s=0.1$，图注写0.2；在获得输入前，该图不能作为精确参数复现的验收样本。

\subsection{陀螺动理学对照：边界、种子和模过滤都要登记}
\citet{E475}的流体实现使用$1024\times128$、双周期伪谱；EUTERPE使用B样条有限元与四阶Runge--Kutta，$x$为Dirichlet而$y$周期，最高约$3\times10^7$标记粒子、最小时间步0.125，$d_e$至少16格。时间归一化须按该文定义换算，不能与前一文的步长直接相比。

流体用确定种子，PIC由噪声开始。边界差别通过局地线性本征模检查受到约束，但非线性比较仍应记录。PIC非线性保留$m=0,\pm1$，流体保留更高谐波；作者认为小岛区高谐波很弱，这不能外推到大岛分裂或双重加速。报告的总能量误差小于2.5\%也应换算成释放磁能比例后，才适合评估加热精度。

\subsection{非局地闭合：刚性与三维模空间必须分开验证}
\citet{L244}用三阶Adams--Bashforth推进，沿周期方向并行。非线性加入热流后，步长可需缩小约100倍；这是所测参数的相对成本，原文没有列出绝对步长和完整网格数量。非局地分母用磁场方向二阶矩的空间平均，分子仍含局地$\nabla_\parallel$，不能改成普通常系数Laplacian而称实现了同一闭合。

线性半盒$(L_x,L_y,L_z)=(46,24\pi,96\pi)$，非线性用$k_y=1/12,k_z=1/48$及$(L_y,L_z)=(12\pi,48\pi)$。单螺旋性允许用$y^*$二维坐标降低成本，却也删去了不同螺旋性的耦合。验证应分别检查闭合算子的线性阻尼、单模增长、能量损失，以及加入多模后的三维演化；任何一个检查不能替代其余检查。
